% §2 Background
\subsection{Cascade Voice Architectures}

A cascade voice architecture chains three independently developed components.

\paragraph{Speech recognition.}
The \emph{automatic speech recognition} (ASR) component converts raw audio input to text.
Modern ASR systems are based on encoder-decoder transformers \citep{radford2022whisper}
or connectionist temporal classification models \citep{graves2006ctc}.
ASR latency has two subcomponents:
\emph{buffering latency}, the time required to accumulate sufficient audio before decoding begins,
which is determined by the end-of-utterance detection (VAD) threshold and cannot be reduced below
the duration of the final word or phrase of the user's turn,
and \emph{decoding latency}, the time to run the recognition model on the buffered audio,
which is reducible with faster hardware or smaller models.
Voice activity detection systems \citep{raux2008optimizing} trigger the end-of-utterance signal
that releases the audio buffer to the ASR model.
The VAD threshold is a critical parameter: a low threshold reduces buffering latency but
increases the rate of premature cut-offs, while a high threshold increases buffering latency
but improves transcription accuracy.
Specaugment \citep{park2019specaugment} and related data augmentation techniques improve
ASR robustness without affecting the latency decomposition.

\paragraph{Language model.}
The language model (LM) component receives the ASR transcript as text input and produces a
text response.
\emph{First-token latency} (also called time-to-first-token, TTFT) is the delay between the LM
receiving input and emitting its first output token.
TTFT is determined by the time to prefill the model's key-value cache with the input context.
For autoregressive transformers \citep{vaswani2017attention}, this is linear in context length
and quadratic in sequence length for self-attention.
Subsequent token generation proceeds at a rate determined by memory bandwidth,
not compute \citep{leviathan2023fast}.
The LM is the only cascade stage that performs open-ended reasoning;
all other stages perform fixed-function transformations.

\paragraph{Speech synthesis.}
The \emph{text-to-speech} (TTS) component converts the LM's text output to audio.
Classic systems such as Tacotron~2 \citep{shen2018natural} and WaveNet \citep{oord2016wavenet}
generate audio non-autoregressively from a complete text input.
Streaming TTS systems can begin audio generation after receiving a short initial text segment,
reducing the time before audio playback begins at the cost of potentially lower quality
for very short prefixes.
The first audio chunk emitted by the TTS model must reach the playback buffer before the
user hears any agent speech;
this buffering delay is part of the irreducible floor when the TTS model cannot begin
synthesis without a minimum text prefix.

\paragraph{Total silence gap.}
The total silence gap in a cascade architecture is the time from the end of the user's utterance
to the beginning of the agent's audio response.
Under the standard (non-streaming) cascade configuration, this equals the sum of:
ASR buffering latency, ASR decoding latency, LM first-token latency, and TTS initial audio latency.
We refer to this sum as the \emph{silence gap} throughout this paper.
The cascade architecture is distinguished from alternative designs precisely by the fact that
these four components contribute additively and sequentially to the silence gap.

\subsection{Native Voice Models}

Native voice models process speech input and produce speech output within a single end-to-end
neural network, eliminating the inter-component boundaries present in cascade architectures.
The key structural difference is that native models can begin producing audio output before
all input audio has been received, enabling sub-200~ms response latencies under favourable
network conditions.

OpenAI's Realtime API \citep{openai2024realtimeapi} exposes a GPT-4o variant that accepts
audio input and produces audio output directly via a WebSocket connection,
with reported response latencies of 320--500~ms.
Google's Gemini Live \citep{google2024geminivoice} provides a similar capability
with streaming audio input and output.
Both systems represent the current state of the art for low-latency voice interaction.

Native voice models eliminate ASR and TTS as separate stages.
They process audio tokens directly, enabling the model to begin generating audio
before the user has finished speaking and to detect barge-in events internally.
The latency of a native voice model is therefore dominated by the time from receiving
sufficient audio context to emitting the first audio output token,
which is bounded below by the model's forward pass time on the input context.

Native voice models support function calling in audio mode~\citep{openai2024realtimeapi};
however, tool-call turns require the model to pause audio generation while the external tool executes,
partially reintroducing sequential latency that increases with the number of tool calls.
This characteristic makes native voice models less suited for applications requiring
real-time access to databases, APIs, or other external services with non-trivial round-trip latency.
Cascade architectures retain a structural advantage for auditable, tool-intensive applications,
motivating the analysis in this paper.

\subsection{Human Turn-Taking and Perceptual Thresholds}

Human conversation is governed by tight timing constraints.
\citet{levinson2015timing} report that the modal inter-turn interval in natural conversation
is approximately 0~ms (overlapping speech) to 200~ms across typologically diverse languages.
\citet{stivers2009universals} confirm this finding across ten unrelated languages,
establishing the 200~ms threshold as a near-universal feature of human turn-taking.
Response gaps exceeding 700--800~ms are perceived as face-threatening or indicative of
uncertainty \citep{levinson2015timing}.
These findings have been replicated in laboratory settings with synthesised speech,
suggesting that the perceptual threshold applies to machine voices as well as human voices.

We adopt $\tau_{human} = 200$~ms as the upper bound on the silence gap within which a voice
agent is perceived as conversationally natural.
We note that this threshold is context-dependent:
task-oriented dialogue (e.g., information retrieval) tolerates longer gaps than casual conversation,
and individual users vary in their sensitivity.
Our formal analysis treats $\tau_{human}$ as a parameter; the 200~ms value is used in worked
examples.
The term \emph{turn-taking threshold} $\tau_{human}$ is used throughout this paper to refer
to this perceptual boundary.

\subsection{Tool-Using Agents and the ReAct Pattern}

The ReAct pattern \citep{yao2023react} interleaves \emph{reasoning} steps, in which the language
model generates a chain-of-thought, with \emph{acting} steps, in which the model calls an
external tool and observes its result.
A $k$-step ReAct agent executes $k$ reasoning--action cycles before producing its final response.
This pattern improves task completion rates on multi-step reasoning benchmarks
but introduces additional latency before the agent can begin its final response.

In voice deployments, ReAct agents are typically implemented using structured output generation:
each LM inference step produces either a \texttt{tool\_call} action (with arguments) or a
\texttt{message} action (the final response to the user).
Frameworks such as LangGraph \citep{langgraph2024} orchestrate these cycles,
routing tool call results back to the LM as observations.
Each cycle adds a full LM inference pass (prefill + decode) and a tool execution round-trip to the
silence gap before the user hears any audio response.
Chain-of-thought reasoning \citep{wei2023chainofthought} improves response quality but
further increases TTFT by extending the generated sequence before the final answer token.
The compounding effect of multiple ReAct cycles on the silence gap is one of the central
analytical results of this paper, developed formally in Section~\ref{sec:react_extension}.
Recent work on voice agent system design \citep{dong2024building} identifies multi-step
tool use as a primary source of perceived latency in deployed cascade voice agents,
motivating the formal treatment we provide here.

\subsection{Spoken Dialogue Systems and Incremental Processing}

The spoken dialogue systems (SDS) community has long studied the sequential latency accumulation problem in cascade speech pipelines~\citep{schlangen2011general,skantze2021turn}.
\citet{schlangen2011general} introduce a general formal model of incremental dialogue processing in which each module produces output incrementally, motivating overlap between pipeline stages as the primary mechanism for latency reduction.
\citet{skantze2021turn} survey turn-taking in conversational systems and human-robot interaction, providing an engineering-facing bridge between psycholinguistic turn-taking norms and deployed dialogue system timing.
\citet{raux2008optimizing} demonstrate empirically how end-of-utterance threshold configuration affects perceived response latency in deployed voice systems.
Our work is distinguished from this literature by deriving a formal lower bound on the irreducible latency floor as a function of pipeline topology, and by connecting this bound explicitly to the human conversational gap tolerance threshold and to multi-step tool-using agents.
